\documentclass[11pt, a4paper]{article}

\usepackage[margin=2.5cm]{geometry}
\usepackage[T1]{fontenc}
\usepackage{lmodern}
\usepackage{microtype}
\usepackage{parskip}
\usepackage{titlesec}
\usepackage{xcolor}
\usepackage{listings}
\usepackage{booktabs}
\usepackage{hyperref}
\usepackage{enumitem}
\usepackage{fancyhdr}
\usepackage{tcolorbox}
\usepackage{graphicx}

% ---- colours ----------------------------------------------------------------
\definecolor{codebg}{HTML}{F5F5F5}
\definecolor{codeframe}{HTML}{CCCCCC}
\definecolor{keywordcol}{HTML}{0060AD}
\definecolor{commentcol}{HTML}{5C8A00}
\definecolor{stringcol}{HTML}{C0392B}
\definecolor{numbercol}{HTML}{7D3C98}
\definecolor{notebg}{HTML}{EAF4FB}
\definecolor{noteframe}{HTML}{5DADE2}
\definecolor{warnbg}{HTML}{FEF9E7}
\definecolor{warnframe}{HTML}{F39C12}
\definecolor{filebg}{HTML}{F0F4F0}
\definecolor{fileframe}{HTML}{3D8B37}

% ---- code listings ----------------------------------------------------------
\lstset{
  language=Python,
  backgroundcolor=\color{codebg},
  frame=single,
  rulecolor=\color{codeframe},
  basicstyle=\ttfamily\small,
  keywordstyle=\color{keywordcol}\bfseries,
  commentstyle=\color{commentcol}\itshape,
  stringstyle=\color{stringcol},
  numberstyle=\tiny\color{numbercol},
  showstringspaces=false,
  breaklines=true,
  breakatwhitespace=true,
  tabsize=4,
  aboveskip=0.8em,
  belowskip=0.8em,
}

% ---- callout boxes ----------------------------------------------------------
\tcbuselibrary{skins, breakable}

\newenvironment{note}{%
  \begin{tcolorbox}[
    colback=notebg, colframe=noteframe,
    fonttitle=\bfseries, title=Note,
    breakable, enhanced
  ]
}{%
  \end{tcolorbox}
}

\newenvironment{intuition}{%
  \begin{tcolorbox}[
    colback=warnbg, colframe=warnframe,
    fonttitle=\bfseries, title=Why does the code do this?,
    breakable, enhanced
  ]
}{%
  \end{tcolorbox}
}

\newenvironment{filebox}[1]{%
  \begin{tcolorbox}[
    colback=filebg, colframe=fileframe,
    fonttitle=\bfseries\ttfamily, title=#1,
    breakable, enhanced
  ]
}{%
  \end{tcolorbox}
}

% ---- section styling --------------------------------------------------------
\titleformat{\section}{\large\bfseries}{}{0em}{}[\titlerule]
\titleformat{\subsection}{\normalsize\bfseries}{}{0em}{}

% ---- page style -------------------------------------------------------------
\pagestyle{fancy}
\fancyhf{}
\lhead{\textit{Codebase Guide -- NeonTreeEvaluation Working Group}}
\rhead{\thepage}
\renewcommand{\headrulewidth}{0.4pt}

% ---- hyperref ---------------------------------------------------------------
\hypersetup{
  colorlinks=true,
  linkcolor=keywordcol,
  urlcolor=keywordcol,
}

% =============================================================================
\begin{document}

\begin{center}
  {\LARGE\bfseries Codebase Guide}\\[0.4em]
  {\large NeonTreeEvaluation Working Group}\\[0.2em]
  {\normalsize\color{gray} A walkthrough of every file and the problem it solves}
\end{center}

\vspace{1em}
\tableofcontents
\vspace{1em}

% =============================================================================
\section{Overview}

The code in this repository does one thing right now: take the raw
NeonTreeEvaluation dataset as downloaded from the internet and transform it
into a form that a PyTorch model can consume directly. None of the model
architecture or training logic is written yet. What exists is the data
pipeline -- the plumbing that sits between the raw files on disk and the
training loop.

The pipeline has to solve several concrete problems that are not obvious until
you actually work with the dataset:

\begin{itemize}[nosep]
  \item The download is large and may fail partway through; we need to detect
        this and avoid silently training on corrupted data.
  \item The archive contains zip files inside zip files; a single unzip is not
        enough.
  \item After unzipping, some directories are double-nested in a way that
        breaks any code that assumes a flat structure.
  \item There are only around 20 source images, each very large. The model
        trains on $400 \times 400$ pixel patches, so we need to extract many
        such patches from each source image and save them individually. This
        also dramatically expands the effective size of the training set.
  \item The raw images are GeoTIFF files. Converting them to PyTorch tensors
        takes several seconds per image. Doing this conversion on every training
        run would be prohibitively slow, so we do it once and save the result.
  \item The RGB imagery and the LiDAR canopy height (CHM) imagery were captured
        at different resolutions, so they do not line up pixel-for-pixel out of
        the box.
  \item The annotation files list bounding boxes in XML format. When a file
        contains only a single tree, the XML parser returns a different data
        structure than when it contains multiple trees, which would silently
        drop single-tree annotations if not handled.
\end{itemize}

The sections below walk through each source file in the order you would
encounter it if you traced a run from \texttt{main.py} downward.

% =============================================================================
\section{\texttt{main.py}}

\begin{filebox}{main.py}
  The entry point. Run this file to execute the full setup pipeline.
\end{filebox}

\begin{lstlisting}
from src.utils.download import download_data, cleanup_files
from src.const import DATA_PATH, NEON_TREE_PATH
from src.data.setup import SetupNeonTreeData
from pathlib import Path

def main():
    if not Path(DATA_PATH).exists():
        download_data(DATA_PATH)
    cleanup_files(NEON_TREE_PATH)
    data = SetupNeonTreeData(force_overwrite=True)
    print("Data loaded successfully.")
\end{lstlisting}

\texttt{main.py} is deliberately thin. It sequences three steps:

\begin{enumerate}[nosep]
  \item Download the raw dataset if the local \texttt{data/} directory does
        not yet exist.
  \item Fix the directory structure problems left by the unzipper
        (see \texttt{download.py}).
  \item Run the full preprocessing pipeline to convert raw images into cached
        PyTorch files.
\end{enumerate}

Having a single entry point like this means anyone on the team can set up the
project from scratch by running one command, without needing to know the
internal details of how the data is obtained or organised.

% =============================================================================
\section{\texttt{src/const.py}}

\begin{filebox}{src/const.py}
  A single place where every project-wide constant lives.
\end{filebox}

\begin{lstlisting}
from pathlib import Path
from src.utils.logger import Logger

DATA_URL    = "https://owncloud.gwdg.de/..."
DATA_PATH   = Path("data/")
DATA_HASH   = "0ad723302caa8f34..."
NEON_TREE_PATH = DATA_PATH / "neon_tree" / "NeonTreeEvaluation"
PT_DATA_PATH   = DATA_PATH / "pt_data"
LOG_DIR     = Path("logs/")
LOGGER      = Logger("neon_tree")

# placeholder values until we compute the real ones from our data
CHANNEL_MEANS  = (0.485, 0.456, 0.406)
CHANNEL_STDEVS = (0.229, 0.224, 0.225)
\end{lstlisting}

The purpose of \texttt{const.py} is to prevent path strings from being
scattered across the codebase. If the dataset ever moves to a different
directory, or we switch hosting providers for the download, there is exactly
one file to edit. Every other file imports its paths from here rather than
hardcoding them.

\texttt{CHANNEL\_MEANS} and \texttt{CHANNEL\_STDEVS} are the normalisation
statistics used to preprocess images before they enter the model. The values
currently in the file are the ImageNet statistics (a standard starting point
from the ML primer), marked with a comment indicating they are placeholders.
The \texttt{normalization\_params.py} utility exists to compute our
dataset-specific values, at which point these constants will be updated.

% =============================================================================
\section{\texttt{src/utils/download.py}}

\begin{filebox}{src/utils/download.py}
  Downloads the dataset, verifies its integrity, and unpacks nested archives.
\end{filebox}

\subsection{Streaming download with a checksum}

The dataset file is large. A naive download that reads the entire response
into memory would crash on any machine without enough RAM, and a download that
saves to disk without verification could silently produce a corrupted file if
the connection dropped partway through.

\texttt{download\_data} solves both problems at once by streaming the response
in 1\,MB chunks, writing each chunk to disk and feeding it into a SHA-256 hash
function simultaneously:

\begin{lstlisting}
h = hashlib.sha256()
for chunk in response.iter_content(chunk_size=1 << 20):
    f.write(chunk)
    h.update(chunk)

digest = h.hexdigest()
if digest != DATA_HASH:
    temp_file.unlink()        # delete the corrupted file
    raise ValueError(...)     # and refuse to continue
\end{lstlisting}

SHA-256 is a mathematical function that produces a fixed-length fingerprint
from any input. The fingerprint stored in \texttt{DATA\_HASH} was computed
from the original, known-good archive. If any byte in the downloaded file
differs -- due to a dropped connection, a server error, or file corruption --
the fingerprint will not match and the code raises an error before any further
processing happens. The incomplete file is deleted automatically so a second
attempt starts fresh rather than skipping the download because a partial file
already exists.

The \texttt{rich} library's \texttt{Progress} class provides the animated
download progress bar. This is purely cosmetic and can be disabled by passing
\texttt{verbose=False}.

\subsection{Recursive unzipping}

The NeonTreeEvaluation archive is a zip file that contains further zip files
inside it (one per data modality). A single \texttt{unzip} call extracts the
outer archive but leaves the inner ones untouched.
\texttt{\_unzip\_recursive} handles this by extracting an archive and then
scanning the result for any further \texttt{.zip} files, repeating until none
remain:

\begin{lstlisting}
def _unzip_recursive(zip_path, dest, verbose=True):
    with ZipFile(zip_path, "r") as zf:
        zf.extractall(dest)

    for extracted in dest.rglob("*.zip"):
        if extracted == zip_path:
            continue
        _unzip_recursive(extracted, extracted.parent / extracted.stem, verbose)
        extracted.unlink()   # delete the .zip once extracted
\end{lstlisting}

Each inner archive is extracted into a subdirectory named after the archive
itself (e.g. \texttt{RGB.zip} extracts into \texttt{RGB/}), and then the
\texttt{.zip} is deleted to save disk space.

\subsection{Fixing double-nested directories}

After extraction, the dataset has a structural quirk: the
\texttt{evaluation/} and \texttt{annotations/} directories each contain a
subdirectory with the same name as themselves. That is, the files end up at
paths like:

\begin{center}
  \texttt{NeonTreeEvaluation/evaluation/evaluation/SJER\_006.tif}
\end{center}

instead of the expected:

\begin{center}
  \texttt{NeonTreeEvaluation/evaluation/SJER\_006.tif}
\end{center}

\texttt{cleanup\_files} detects this pattern and collapses it by moving every
item from the inner directory up to the outer one:

\begin{lstlisting}
def remove_nested(save_path, target_name):
    dir        = save_path / target_name
    nested_dir = dir / target_name          # e.g. evaluation/evaluation
    if nested_dir.exists():
        for item in nested_dir.iterdir():
            shutil.move(str(item), dir / item.name)
        nested_dir.rmdir()

remove_nested(save_path, "evaluation")
remove_nested(save_path, "annotations")
\end{lstlisting}

Without this step, every file-loading function would need to know about the
double nesting and construct paths accordingly. Fixing it once at setup time
means the rest of the code can assume the straightforward structure.

% =============================================================================
\section{\texttt{src/utils/singleton.py}}

\begin{filebox}{src/utils/singleton.py}
  A metaclass that ensures a class can only ever be instantiated once.
\end{filebox}

\begin{lstlisting}
class SingletonMeta(type):
    _instances = {}

    def __call__(cls, *args, **kwargs):
        if cls not in cls._instances:
            instance = super().__call__(*args, **kwargs)
            cls._instances[cls] = instance
        return cls._instances[cls]
\end{lstlisting}

A \emph{metaclass} is a class whose instances are themselves classes. That is
a confusing sentence, so here is the practical effect: any class that uses
\texttt{SingletonMeta} as its metaclass will, the second time you try to
create an instance, just hand you back the instance that was created the first
time.

\begin{intuition}
  \texttt{SetupNeonTreeData} (described in Section~\ref{sec:setup}) is a
  singleton because the preprocessing work it does -- reading every image from
  disk and saving it as a \texttt{.pt} file -- is expensive and should only
  run once per execution. If multiple parts of the codebase each tried to
  construct a \texttt{SetupNeonTreeData} object independently, without the
  singleton guarantee, they would each kick off the full preprocessing run
  again. The singleton ensures that only the first construction does real work;
  subsequent constructions silently return the already-initialised object.
\end{intuition}

% =============================================================================
\section{\texttt{src/utils/logger.py}}

\begin{filebox}{src/utils/logger.py}
  A logging wrapper that writes structured records to a file and pretty output
  to the terminal simultaneously.
\end{filebox}

\begin{lstlisting}
class Logger:
    def __init__(self, name):
        self.console = Console()          # rich terminal output
        self._logger = self._setup_logger(name)

    def log_and_print(self, rich_object):
        # render the object to plain text and write to the log file
        ...
        self._logger.info(plain_text)
        # also render it to the terminal with colour and formatting
        self.console.print(rich_object)

    def __getattr__(self, name):
        return getattr(self._logger, name)   # delegate .debug(), .warning() etc.
\end{lstlisting}

Python's built-in \texttt{logging} module writes plain text to log files.
The \texttt{rich} library renders coloured, formatted output to the terminal.
\texttt{Logger} combines both: \texttt{log\_and\_print} sends output to both
destinations at once, while the \texttt{\_\_getattr\_\_} delegation means you
can also call \texttt{LOGGER.debug(...)}, \texttt{LOGGER.warning(...)} and so
on exactly as you would with a standard Python logger (these go to the file
only, without terminal output).

Each run of the program creates a new log file named with the current
timestamp (e.g. \texttt{logs/2026-03-18\_14-32-07.log}). This means you can
look back at the output from previous runs without it being overwritten.

A single shared \texttt{LOGGER} instance is created in \texttt{const.py} and
imported wherever it is needed, which keeps all log records for a given run in
one file.

% =============================================================================
\section{\texttt{src/utils/transforms.py}}

\begin{filebox}{src/utils/transforms.py}
  A custom Torchvision transform that converts a raw image array into a
  normalised float tensor.
\end{filebox}

\begin{lstlisting}
class ToTensor(torch.nn.Module):
    def forward(self, img, label=None):
        converted = v2.Compose([
            v2.ToImage(),
            v2.ToDtype(torch.float32, scale=True),
        ])(img)

        if label is not None:
            return converted, label
        else:
            return converted
\end{lstlisting}

Two conversions happen here in sequence:

\begin{enumerate}[nosep]
  \item \texttt{v2.ToImage()} converts the NumPy array returned by the image
        reader into a PyTorch tensor and reorders the dimensions from
        \texttt{(H, W, C)} to \texttt{(C, H, W)} -- PyTorch's expected
        convention.
  \item \texttt{v2.ToDtype(torch.float32, scale=True)} changes the numeric
        type from unsigned 8-bit integers (values 0--255) to 32-bit floats
        (values 0.0--1.0). The \texttt{scale=True} argument performs the
        division by 255 automatically.
\end{enumerate}

\begin{intuition}
  Why divide by 255? Neural network optimisers (the algorithms that update
  the model's weights during training) are sensitive to the scale of the
  numbers they see. Input values of 0--255 and weight values of approximately
  0.0--1.0 are on very different scales, which makes training unstable.
  Rescaling the inputs to 0.0--1.0 keeps everything in a compatible range.
\end{intuition}

The transform accepts an optional \texttt{label} argument and passes it
through unchanged. This allows it to be used in contexts where an image and
its annotation are processed together without accidentally discarding the
annotation.

% =============================================================================
\section{\texttt{src/utils/normalization\_params.py}}

\begin{filebox}{src/utils/normalization\_params.py}
  A utility script for computing the channel-wise mean and standard deviation
  of the dataset, to be run once and used to update \texttt{CHANNEL\_MEANS} and
  \texttt{CHANNEL\_STDEVS} in \texttt{const.py}.
\end{filebox}

\begin{lstlisting}
def get_images(data_dir):
    files     = sorted(data_dir.glob("*.pt"))
    unpickles = [torch.load(file, weights_only=False) for file in files]
    images    = [np.asarray(image) for image, boxes in unpickles]
    return np.stack(images)
\end{lstlisting}

This file is a standalone script (note the \texttt{if \_\_name\_\_ == "\_\_main\_\_"} block at the bottom) rather than a module imported by the rest of
the code. Its job is to iterate over the cached \texttt{.pt} files, load the
image tensors, stack them into a single large array, and compute statistics
across the whole training set.

\begin{note}
  This script is not yet complete. The \texttt{means()} function currently
  computes the per-channel means but does not return them -- a missing
  \texttt{return} statement -- and there is no corresponding function for the
  standard deviations. Once these are implemented, the values should replace
  the ImageNet placeholders in \texttt{const.py}.

  The reason we compute these from our own data rather than reusing ImageNet
  values permanently is that our images are aerial photographs, which have a
  different brightness distribution than ground-level photographs.
  Dataset-specific normalisation values will give the model a more consistent
  input distribution to work with.
\end{note}

% =============================================================================
\section{\texttt{src/data/setup.py}}
\label{sec:setup}

\begin{filebox}{src/data/setup.py}
  The main preprocessing class. Reads raw TIFFs and XML annotations, combines
  RGB and CHM channels, extracts $400 \times 400$ crops, and saves each crop
  as a \texttt{.pt} file for fast loading during training.
\end{filebox}

\subsection{Why preprocess and cache?}

Loading a raw GeoTIFF, resizing it, converting it to a tensor, and parsing an
XML file takes several seconds per image. A training run with thousands of
iterations, each loading a new image, would spend more time on disk I/O than
on actual learning. The standard solution is to do all the slow preprocessing
once and save the result in a format that is fast to load -- in our case,
PyTorch's native \texttt{.pt} format.

\texttt{SetupNeonTreeData} uses the singleton pattern from
\texttt{singleton.py} to guarantee that this preprocessing only ever runs once
per program execution, regardless of how many times different parts of the code
ask for a \texttt{SetupNeonTreeData} object.

\subsection{The preprocessing loop}

The dataset contains around 20 source images, each capturing a large plot area
at high resolution. These are far bigger than the $400 \times 400$ pixel
patches the model will train on. Rather than loading a huge image and cropping
at training time (which wastes I/O bandwidth re-reading the same source
repeatedly), setup extracts all crops up front and saves each one as its own
\texttt{.pt} file. This also means the number of training examples is no
longer limited by the number of source images -- we can extract as many
non-overlapping (or overlapping) patches as the source images allow.

For each RGB image in the training or evaluation split, the code:

\begin{enumerate}
  \item Finds the corresponding annotation file (in \texttt{annotations/}) and
        CHM file (in \texttt{CHM/}) by constructing the expected paths from the
        image filename.
  \item Skips the image with a warning if either the annotation or the CHM is
        missing -- a graceful fallback rather than a crash.
  \item Loads the RGB image as a tensor of shape \texttt{(3, H, W)}.
  \item Loads the CHM image and resizes it to match the RGB spatial dimensions,
        since the two modalities were captured at different resolutions.
  \item Concatenates RGB and CHM along the channel dimension, producing a
        4-channel tensor of shape \texttt{(4, H, W)}.
  \item Parses the bounding boxes from the XML annotation file.
  \item Extracts $400 \times 400$ crops from the full image. For each crop, the
        bounding boxes are filtered and clipped to the crop window: boxes that
        fall entirely outside the window are discarded, and boxes that overlap
        the edge are clipped to the boundary.
  \item Saves each \texttt{(crop\_tensor, cropped\_bounding\_boxes)} pair as a
        separate \texttt{.pt} file, named by source image and crop index.
\end{enumerate}

\begin{intuition}
  Doing the crop-and-clip step at setup time rather than training time is what
  makes this approach work. At training time, the \texttt{Dataset} class simply
  loads a file and gets back a correctly sized image with correctly adjusted
  annotations, with no further geometric computation needed. The expensive work
  -- file I/O, coordinate arithmetic, box filtering -- happens exactly once.
\end{intuition}

\subsection{Merging RGB and CHM}

The channel concatenation step is the most consequential data design decision
in the codebase so far:

\begin{lstlisting}
rgb  = self._load_image(img_path)                    # shape: (3, H, W)
chm  = self._load_image(chm_path, target_size=rgb.shape[1:])  # shape: (1, H, W)
comb = torch.cat([rgb, chm], dim=0)                  # shape: (4, H, W)
\end{lstlisting}

By stacking LiDAR height as a fourth channel alongside the three colour
channels, the model can learn to use height information alongside colour when
deciding whether a blob in the image is a tree crown. This is one of the main
reasons to use the NEON dataset over ordinary aerial photography -- the height
channel makes it much easier to distinguish tall trees from low shrubs that
might look similar in colour alone.

\subsection{Parsing bounding box annotations}

The annotation files are in Pascal VOC XML format, the same format used in
many standard computer vision datasets. A typical file looks like:

\begin{lstlisting}[language=XML]
<annotation>
  <object>
    <bndbox>
      <xmin>45</xmin> <ymin>112</ymin>
      <xmax>130</xmax> <ymax>198</ymax>
    </bndbox>
  </object>
  <object> ... </object>
</annotation>
\end{lstlisting}

The \texttt{xmltodict} library parses this into a Python dictionary. There is
a subtle issue: when the file contains multiple \texttt{<object>} entries,
\texttt{xmltodict} puts them in a list. When there is exactly one entry, it
puts it in a dictionary directly (not inside a list). The code handles both
cases explicitly:

\begin{lstlisting}
if isinstance(data["annotation"]["object"], list):
    bounding_boxes = [[...] for object_dict in data["annotation"]["object"]]
elif isinstance(data["annotation"]["object"], dict):
    bounding_boxes = [[...]]   # wrap the single dict in a list
\end{lstlisting}

Without this check, images containing only one tree would silently fail to
produce any annotations, corrupting the training data for those samples.

The extracted coordinates are wrapped in a \texttt{BoundingBoxes} object from
Torchvision, which stores not just the coordinates but also the format
(\texttt{"XYXY"} -- two corner points, as opposed to centre-plus-size) and
the canvas size (the image dimensions). Storing bounding boxes in this typed
wrapper rather than as a plain tensor means that Torchvision's crop and flip
transforms -- applied both during setup and at training time -- automatically
update the coordinates to match whatever geometric operation was applied to the
image. If the boxes were stored as a plain \texttt{torch.Tensor}, a crop would
move the pixels but leave the coordinates pointing at the wrong locations.

% =============================================================================
\section{\texttt{src/dataset.py}}

\begin{filebox}{src/dataset.py}
  The PyTorch \texttt{Dataset} class that the training loop will consume.
\end{filebox}

\begin{lstlisting}
class TreeImageDataset(Dataset):
    def __init__(self, input_dir, transforms=None):
        SetupNeonTreeData()   # ensure preprocessing has been run
        self.paths = sorted(input_dir.glob("*.pt"))
        self.transform = Compose(transforms) if transforms else Compose([
            Normalize(CHANNEL_MEANS, CHANNEL_STDEVS),
            RandomHorizontalFlip(),
            RandomVerticalFlip(),
            SanitizeBoundingBoxes(),
        ])

    def __len__(self):
        return len(self.paths)

    def __getitem__(self, idx):
        file, boxes = torch.load(self.paths[idx])
        file = self.transform(file)
        return file, boxes
\end{lstlisting}

\texttt{TreeImageDataset} is the class that the \texttt{DataLoader} described
in the ML primer will wrap. It satisfies PyTorch's Dataset contract by
implementing \texttt{\_\_len\_\_} (how many examples?) and
\texttt{\_\_getitem\_\_} (give me example number \texttt{idx}).

Because \texttt{setup.py} has already extracted $400 \times 400$ crops and
saved each one as its own \texttt{.pt} file, \texttt{\_\_getitem\_\_} is
fast: one \texttt{torch.load} call returns an image that is already the right
size with already-adjusted annotations. No image decoding, no XML parsing, no
resizing, no cropping -- all of that happened once during setup.

The default transform pipeline does the following to each example before
returning it:

\begin{enumerate}[nosep]
  \item \textbf{Normalize} -- subtracts the channel means and divides by
        the standard deviations, shifting the pixel value distribution to be
        centred near zero. This is expected by the pre-trained model
        architecture.
  \item \textbf{RandomHorizontalFlip / RandomVerticalFlip} -- randomly mirrors
        the image left-right and/or top-bottom. Aerial imagery has no canonical
        orientation (north can point in any direction), so all four orientations
        are equally valid and showing the model all of them improves robustness.
  \item \textbf{SanitizeBoundingBoxes} -- cleans up any bounding boxes whose
        coordinates became degenerate after a flip. This is mainly a safety
        measure; correctly formed boxes should survive a flip unchanged in
        terms of validity.
\end{enumerate}

Note that there is no \texttt{RandomCrop} in this pipeline. That step has been
moved entirely into \texttt{setup.py}. This is a deliberate trade-off: we lose
the ability to see a different random crop of each source image on every
training epoch (since the crops are now fixed), but we gain a much faster and
simpler training loop. If we want more variety, we can simply extract more
crops during setup -- overlapping ones, or crops taken at different random
offsets -- without touching the training code at all.

\begin{note}
  The constructor calls \texttt{SetupNeonTreeData()} without arguments at
  startup. Because \texttt{SetupNeonTreeData} is a singleton, this either
  triggers preprocessing (if it has not been done yet) or returns the existing
  instance immediately (if it has). This means you can construct a
  \texttt{TreeImageDataset} without having run \texttt{main.py} first; the
  dataset will set itself up automatically.
\end{note}

% =============================================================================
\section{How the Files Fit Together}

\begin{center}
  \begin{tabular}{lll}
    \toprule
    \textbf{File}                           & \textbf{Role}      & \textbf{Key output}           \\
    \midrule
    \texttt{main.py}                        & entry point        & runs the full setup           \\
    \texttt{const.py}                       & constants          & paths, hash, logger, stats    \\
    \texttt{utils/download.py}              & fetch data         & verified, flat directory tree \\
    \texttt{utils/singleton.py}             & pattern            & one-instance guarantee        \\
    \texttt{utils/logger.py}                & observability      & file + terminal logging       \\
    \texttt{utils/transforms.py}            & conversion         & float tensor from raw array   \\
    \texttt{utils/normalization\_params.py} & statistics         & dataset mean/stdev (WIP)      \\
    \texttt{data/setup.py}                  & preprocessing      & \texttt{.pt} files on disk    \\
    \texttt{dataset.py}                     & training interface & batches of (image, boxes)     \\
    \bottomrule
  \end{tabular}
\end{center}

The flow is linear: \texttt{main.py} calls \texttt{download.py} to fetch and
flatten the raw data, then calls \texttt{setup.py} to convert it into
\texttt{.pt} files. From that point on, \texttt{dataset.py} is the only file
that the training loop needs to interact with -- everything upstream of it has
already done its work.

\end{document}
